I would like to express my sincere gratitude to Professor Didem Gürdür Broo for providing the opportunity to undertake this thesis within the Cyber-Physical Systems Lab (CPS Lab) at Uppsala University, and to my supervisor, Xuezhi Niu, for his guidance, constructive feedback, and support throughout this research.

I am grateful to my examiner, Ramanathan Thinniyam Srinivasan, for his valuable comments and suggestions, which helped improve the quality of this thesis.

This thesis has been a challenging and rewarding experience. It required bringing together knowledge from multiple disciplines, including warehouse automation, heterogeneous multi-robot systems, and large language models. Combining these areas into a single research project demanded continuous learning and a willingness to explore unfamiliar domains.

The rapid development of large language models made this research particularly challenging. Keeping pace with the latest developments while integrating them into a grounded decision-making framework and evaluating them through extensive experiments required continuous learning, patience, and significant computational resources.

Throughout this journey, I would like to thank my family for their unwavering love, support, encouragement, and belief in me, which gave me the strength to keep going.

Finally, I would like to acknowledge the developers and maintainers of the open-source software, libraries, and simulation environments that made this research possible. Their contributions have been invaluable in enabling research such as this and continue to advance the broader scientific community.
